\section{Conclusion}


We have demonstrated that post-quantum cryptography is viable on current EVM infrastructure through careful optimization of the Number Theoretic Transform. Our 85\% gas reduction—from 24M to 3.6M for FALCON-512 verification—brings quantum-resistant signatures within practical reach for production blockchain systems.

This work represents a critical step toward quantum-safe blockchain infrastructure. By providing efficient NTT operations as a foundational primitive, we enable the deployment of NIST-standardized post-quantum algorithms including Dilithium, FALCON, and SPHINCS+ on EVM-compatible chains. The proposed EIP for NTT precompiles would benefit the entire Ethereum ecosystem, providing a 10-100× improvement over smart contract implementations.

As quantum computing advances rapidly, the blockchain community must act decisively to ensure long-term security. Our open-source implementation, available at \url{https://github.com/lux-network/ntt-evm}, provides an immediate path forward for projects seeking quantum resistance. We encourage the Ethereum community to consider our EIP proposal and begin the transition to post-quantum cryptography before it becomes critical.

The journey from 24M to 3.6M gas demonstrates that with careful optimization, even computationally intensive cryptographic operations can be made practical on constrained environments like the EVM. As the first production-ready NTT implementation for EVM, this work paves the way for a quantum-resistant future while maintaining full backward compatibility with existing smart contract ecosystems.

